\section{Related Work}
\label{sec:related_work}

\noindent \textbf{Optimization-based NVS.}
Since the introduction of Neural Radiance Fields (NeRF)~\cite{mildenhall2020nerf}, the research on novel view synthesis (NVS) has made remarkable progress, with numerous extensions~\cite{yu2021plenoctrees, fridovich2022plenoxels, fastnerf, efficient_nerf, muller2022instant, tensorf, cao2023hexplane, fridovich2023k, bui2025moblurf}. 3D Gaussian Splatting (3DGS)~\cite{kerbl20233d} has emerged as an explicit scene representation approach that models a scene using millions of anisotropic Gaussian primitives and employs differentiable rasterization to achieve photorealistic novel views with real-time inference. Despite these advantages, 3DGS relies on per-scene optimization of Gaussians and cannot directly render novel views from input images without retraining. Very recently, several post-optimization techniques have been proposed to control the number of Gaussian primitives~\cite{HansonTuPUP3DGS, fan2023lightgaussian}. 
PUP 3D-GS~\cite{HansonTuPUP3DGS} computes a spatial sensitivity score to prune pretrained 3DGS models.
LightGaussian~\cite{fan2023lightgaussian} utilizes significance-based pruning, SH coefficient distillation, and vector quantization to reduce the 3DGS model size.
However, similar to 3DGS~\cite{kerbl20233d}, these methods cannot operate in a feed-forward manner.

\begin{figure*}
\centering
\includegraphics[width=\linewidth,keepaspectratio]{figure/EcoSplat_architecture.pdf}
\vspace{-0.6cm}
\caption{\textbf{Overview of EcoSplat.} EcoSplat is trained in two stages: Pixel-aligned Gaussian Training (PGT) (Sec.~\ref{sec:architecture}) and Importance-aware Gaussian Finetuning (IGF) (Sec.~\ref{sec:IGF}).
During IGF, the combination of the importance-aware opacity loss $\mathcal{L}_\text{io}$ and the Progressive Learning on Gaussian Compaction (PLGC) encourages EcoSplat to suppress the opacities of less important Gaussians. At inference, it adaptively satisfies an arbitrary user-specified primitive count and produces the optimal Gaussians in a feed-forward manner (Sec.~\ref{sec:inference}).}
\label{fig:overall_architecture}
\end{figure*}

\noindent \textbf{Feed-forward NVS.} 
To enable generalizable NVS, early methods learn generalizable neural radiance fields~\cite{pixelnerf, mvsnerf, nerf_attention} or image-based rendering models~\cite{ibrnet, neuray, gpnr}, leveraging data-driven priors via transformer-based architectures. More recently, due to the superior rendering quality and efficiency of 3DGS, several feed-forward approaches~\cite{charatan2024pixelsplat, szymanowicz2024splatter, chen2024mvsplat, xu2025depthsplat, kang2025selfsplat, ye2024no, huang2025no, zhang2024gs} have been proposed to directly generate 3D Gaussians from input images. These methods predict the parameters of a single 3D Gaussian per pixel in each input image to represent the target scene. PixelSplat~\cite{charatan2024pixelsplat} and Splatter Image~\cite{szymanowicz2024splatter} are pioneering works that predict 3D Gaussians directly from image features. MVSplat~\cite{chen2024mvsplat} constructs a cost volume via plane sweeping to capture cross-view feature correspondences, while DepthSplat~\cite{xu2025depthsplat} leverages pre-trained monocular depth features~\cite{yang2024depth} to achieve more robust 3D Gaussian reconstructions. Several methods~\cite{hong2024unifying, kang2025selfsplat, li2024ggrt, ye2024no, huang2025no} have been studied for pose-free feed-forward NVS. However, all of the aforementioned feed-forward methods are constrained by their \textit{pixel-aligned Gaussian} representations, which inevitably introduce redundant Gaussians, limiting their efficiency and scalability for dense multi-view scenarios.

\noindent \textbf{Efficient Feed-forward NVS.} 
To address the inefficiency of the above \textit{pixel-aligned} feed-forward NVS methods, recent studies~\cite{jiang2025anysplat, zhang2024gaussian, ziwen2025long, liu2025worldmirror} have proposed strategies for Gaussian aggregation and pruning. AnySplat~\cite{jiang2025anysplat} and WorldMirror~\cite{liu2025worldmirror}, built upon the VGGT~\cite{wang2025vggt} architecture, cluster Gaussians into fixed-size voxels via differentiable voxelization, but the voxel size is a sensitive hyperparameter that strongly affects rendering quality. GGN~\cite{zhang2024gaussian} models the inter-view relationships among Gaussian groups using a graph-based approach and prunes redundant Gaussians via pooling. However, this graph-based approach has been shown to lack sufficient robustness, which hinders its generalization across diverse scene geometries and dense multi-view inputs.
Long-LRM~\cite{ziwen2025long} utilizes a Mamba2~\cite{dao2024transformers}-based architecture to compress long sequences of Gaussians through token merging and opacity-based pruning. Despite of their advancements in efficient rendering, the aforementioned approaches provide neither explicit control over the desired number of primitives nor a guaranteed, optimal trade-off between rendering fidelity and primitive count.
